\section{Operational detail: sources, operators, and retained state}
\label{sec:operations}
The preceding sections describe the engine's components. This section follows
their execution boundaries: when a source is read, what an operator retains
between pulls, and when output becomes available. These concrete mechanisms
describe the host runtime; the compiled route must provide equivalent
functionality under the contract in Section~2.

\subsection{Binding a source and deciding what to read}
A registered materialized table already owns or shares its column buffers.
A lazy table instead starts with a schema, a row count, and a decoder supplied
by its source. This allows a caller to resolve column demand before asking for
column values. \code{LazyTable::project} finds demanded columns missing from
its cache, decodes them, and returns columns in source-schema order. Its cache
contains whole-source columns, so subsequent uses can reuse those buffers.

Filtering adds a second decision: which rows need payload values? In the
ordinary \code{project_where} path, the engine obtains columns needed by
predicates, evaluates the predicates, and constructs a selection of surviving
source-row indices. It then requests the demanded payload columns for those
rows. Already cached columns can be gathered in memory. Newly decoded selected
payloads are not inserted as whole-source cache entries: otherwise a later
query could mistake a filtered subset for the original column.

For example, a scan that filters ticks on \code{price} and returns only
\code{ts} and \code{volume} needs the predicate's answer but does not need
\code{price} in its output. An eligible reader can evaluate certain range or
string predicates while decoding, returning the selection without constructing
the full predicate column. If the reader reports that it cannot answer that
request, the engine uses ordinary decode-and-filter execution. An error is
different from a declined optimization and propagates to the caller.

This separates three quantities: columns referenced by predicates, columns
required by consumers, and source rows selected for those consumers. It also
explains why column-demand analysis and decoding cannot be understood solely
by looking at the final projection.

\evidence{\src{src/runtime/lazy_table.cpp}, \code{LazyTable::project},
\code{project_where}, \code{decode_whole_columns}, and \code{project_rows};
\src{include/ibex/runtime/lazy_table.hpp}, \code{ColumnDecodeFn} and
\code{LazySourceReader}.}

\subsection{Source units and reader ownership}
A lazy source may expose a sequence of \code{SourceUnit} ranges that covers
its rows in order. A reader without such a decomposition is read as a whole.
For a decomposed source, a unit restricts decoding to part of the source;
selection indices still refer to the original source, not to positions
starting at zero within that unit. The reader intersects the selection with
the unit. Keeping one index space prevents a selection made before decoding
from selecting different rows after partitioning.

Each \code{LazySourceReader} owns its mutable decoder state. The reader
factory and pool permit independent readers to share immutable resources
without sharing a decoder cursor. In a scan pipeline, workers can produce
different source units concurrently, while the sequence-indexed output buffer
controls when the consumer receives them. Buffer capacity applies
backpressure: production must wait when the consumer has not released space.
The unit size belongs to the source decomposition; it is distinct from the
morsel grain used to partition an already materialized input.

\evidence{\src{include/ibex/runtime/lazy_table.hpp},
\code{LazySourceReader::decode_units} and \code{decode};
\src{src/runtime/lazy_table.cpp}, \code{acquire_reader} and
\code{release_reader}; \src{src/runtime/pipeline_executor.cpp},
\code{PipelinedScanOperator}.}

\subsection{Map operations: sharing, evaluating, and gathering}
A projection or rename can change a chunk's column entries without visiting
its value rows. The \code{map_chunk} kernel receives source column positions
and output names and constructs entries sharing the selected column values.
The caller supplies the derived row-layout properties. This is a metadata
operation over the columns, rather than a copy of every selected value.

A filter changes the row set and therefore has different work to do. The
predicate identifies survivors; gathering constructs output columns for those
positions, together with their validity. An update evaluates expressions to
produce replacement or additional columns. A fused filter/project step can
gather only the columns required by its output. The physical map description
records kernel capability and column representation/null-policy signatures,
so representation-specific work is chosen at the execution boundary.

Fixed-width kernels can borrow \code{ColumnView} objects containing a value
pointer, a length, and a validity pointer. These views own no storage: the
input chunk or table must remain alive during their use. A string or packed
boolean column requires its own access pattern, rather than being treated
as an ordinary array of independently addressable scalar objects.

\evidence{\src{src/runtime/kernel_types.hpp}, \code{ColumnView},
\code{ChunkView}, and \code{map_chunk};
\src{include/ibex/runtime/pipeline.hpp}, \code{MapStep},
\code{MapKernelCapability}, and \code{ColumnKernelSignature};
\src{tests/test_physical_plan.cpp}, projection and fused-filter tests.}

\subsection{Hash joins: build, probe, and output assembly}
For an eligible equality join, the build phase constructs a lookup structure
over one input. The probe phase looks up keys from the other input and gathers
the matching payload columns. These phases have different lifetime needs:
the index and its backing table must survive every probe, while probe chunks
can be processed one at a time.

\paragraph{Choosing the build side.}
In the ordinary single-key streaming path, the right input is already
materialized. If it is sufficiently small, the engine builds its index without
draining the left input. Otherwise it buffers left chunks until their row count
reaches the right-side count or the left input ends. If the count reaches the
right-side count, it builds on the right and replays the buffered chunks before
continuing the left stream. Only the exhausted smaller-left case needs to
concatenate those chunks to consider building on the left. An ordering benefit
can still favor building on the right. Thus build-side selection depends on
both input size and useful ordering; it is not simply ``always hash the
smaller table''.

\paragraph{Representing repeated keys.}
\code{JoinHashIndex} stores a head row for each key and a \code{chain_next}
array linking build rows with the same key. A uniqueness flag identifies the
case where chain traversal can be avoided. In the serial build, reverse row
insertion makes duplicate chains traverse in original build-row order.
Null keys are excluded in this streaming path, which implements the
\code{nulls never} case. Other join semantics select other implementations.
For categorical probes, the engine can resolve dictionary strings to index
heads once per dictionary and then use row codes to access those heads.

\paragraph{Producing rows.}
Probing produces pairs of left and right row indices. Output assembly gathers
the requested payloads through those pairs and applies the planned output
names. Repeated keys can expand one probe row into many result rows, so memory
for matches depends on result cardinality as well as input cardinality.
Worker-private probe state retains shared ownership of the immutable build.
The orientation and ordering logic must be considered when describing output
order; the order of hash buckets is not a semantic ordering rule.

\paragraph{Filtering a deferred probe source.}
A separate deferred path obtains build keys before decoding the probe payload.
It publishes a dynamic key filter and marks the filter ready. A supported
source can use that filter to narrow candidate rows; exact join probing still
determines matches. The two-phase path can find matching key rows first and
materialize payload columns for survivors afterwards. If that optimization
declines, execution returns to the ordinary materialization/probe path. This
is why a deferred probe source cannot always be pulled independently before
the join's build phase has run.

\evidence{\src{src/runtime/join_chunked.cpp},
\code{ChunkedInnerJoinOperator::initialize}, \code{BufferedThenStreamSource},
\code{choose_and_build_single_key}, \code{JoinHashIndex},
\code{fill_partitioned_heads}, \code{JoinProbe::resolve_categorical_heads},
\code{JoinProbe::assemble_output}, \code{resolve_deferred_probe}, and
\code{try_two_phase_probe}; \src{src/runtime/join_chunked_internal.hpp},
\code{JoinProbeFactory}. This account covers the inspected streaming equality
join family, not every join kind.}

\subsection{Aggregation: state survives the input chunks}
Aggregation replaces many input rows with group results. In the hash path,
group state survives across chunk pulls, while each input chunk is released
after its contribution has been accumulated. The coordinator enforces four
phases:
\begin{enumerate}
  \item \emph{Discovery} identifies the group for each input row and retains
    the active chunk while its column references are in use.
  \item \emph{Accumulation} consumes the discovered group identifiers and
    updates aggregate state, then releases the chunk. Some specialized paths
    combine discovery and accumulation; an explicit transfer marker prevents
    accumulating those values twice.
  \item \emph{Final ordering}, after input exhaustion, brings partition-owned
    groups into global first-occurrence order. Strategies that already use
    global group identifiers need no deferred ordering merge.
  \item \emph{Emission} constructs result columns from finalized group state.
    The implementation rejects emission before input consumption or ordering
    finalization is complete.
\end{enumerate}
For a grouped sum, the conceptual state is an accumulator for each group;
processing row $i$ updates the accumulator selected by its group identifier.
The implementation specializes discovery and accumulation by key and value
representation. Distinct counting needs additional per-group value state,
so ``one accumulator per group'' is not a memory model for every aggregate.
The hash path can consume chunks without retaining all their original rows,
but it emits its result only after the input has ended.

\paragraph{When sorted input changes the algorithm.}
The adaptive aggregate examines the first nonempty chunk and its ordering
metadata. The sorted path requires the leading ordering keys to cover the
grouping columns and rejects grouping entries carrying validity bitmaps.
There are also aggregate-specific gates: distinct counting and nonnumeric
\code{first}/\code{last}, for example, use the hash fallback in this path.
On fallback, the inspected first chunk is prepended to the remaining stream;
it is neither discarded nor read again from the source.

When the sorted strategy is admitted, equal keys occupy contiguous runs.
The operator retains the current group's state across chunk boundaries,
closes it when the key changes, and buffers completed groups for output.
At end of input it closes the remaining open group. The relevant ordering
is ordering by the grouping keys: a time-ascending table with interleaved
symbols does not automatically qualify for aggregation by symbol. This is
an operational reason to distinguish a TimeFrame's time order from a
group-contiguous layout.

\evidence{\src{src/runtime/aggregate_chunked.cpp},
\code{HashAggregateState::next_discovery}, \code{accumulate_discovery},
\code{finalize_ordering}, \code{emit_output},
\code{HashAggregatePhaseOperator::next}, and
\code{ChunkedSortedAggregateOperator}, especially \code{decide_strategy},
\code{sorted_on_group_by}, \code{needs_hash_fallback}, and \code{next_sorted}.}

\subsection{Interchange and memory lifetime}
The Arrow C Data Interface supplies array/schema descriptors and a release
callback for sharing columnar data within a process~\cite{arrowCData}.
Ibex's borrowing import copies values. Its adopting import instead creates a
shared owner for the producer's array and, after successful validation and
import, clears the caller's array descriptor. The retained owner's destructor
eventually calls the producer's release function. Failed adoption leaves
ownership with the caller.

Whether a payload is adopted or converted depends on representation. The
inspected implementation can borrow supported primitive buffers and validated
UTF-8 buffers through that owner. Narrow integer inputs are widened into
Ibex's integer representation; timestamps requiring a unit conversion also
need new value storage. Sharing an external buffer therefore preserves a
producer lifetime, while converting it allocates an Ibex representation.
The implementation's \code{import_column} dispatch is the evidence here:
the public header's shorter adoption comment lists fewer supported layouts.

\evidence{\src{src/interop/arrow_c_data.cpp}, \code{AdoptedArrayOwner},
\code{import_column}, \code{import_string_column},
\code{import_table_from_arrow}, and \code{adopt_table_from_arrow}.}

For operational reasoning, the following lists retained state rather than
asserting an unmeasured whole-query memory bound. Shared buffers must be
counted once even when several entries refer to them.
\begin{center}
\begin{tabularx}{\linewidth}{@{}lX@{}}
\toprule
Boundary & Storage that may remain live\\
\midrule
Lazy binding & Whole-source columns already decoded into its cache.\\
Map pipeline & Input buffers, worker scratch/output, and buffered chunks.\\
Hash join & Build table and index; buffered probe input and match indices.\\
Hash aggregate & Group keys and aggregate state across all input chunks.\\
Sorted aggregate & Open-group state and buffered completed-group output.\\
Result sink & The growing result and the current incoming chunk.\\
Arrow adoption & Producer storage retained by the shared array owner.\\
\bottomrule
\end{tabularx}
\end{center}
The distinction between input consumption and output production is central:
a bounded chunk interface can feed an operator whose retained state grows
with the number of groups, join matches, or result rows.

\section{A trade-summary workload}
\label{sec:trade-workload}
The runnable example \src{docs/engine/trade_summary.ibex} starts with five
timestamped ticks in ascending time order. It constructs a TimeFrame, removes
a zero-price tick, computes trade notionals, enriches symbols with an instrument
table, and produces volume-weighted average prices (VWAP). This first workload
uses a complete sample interval; window and resampling rules remain deferred.

\lstinputlisting[firstline=15]{trade_summary.ibex}

The instrument table has one row per symbol, with venue \code{XNAS} for both
symbols. The expected data flow is:
\begin{center}
\begin{tabularx}{\linewidth}{@{}lrlX@{}}
\toprule
Stage & Rows & Operation & Resulting work\\
\midrule
Ticks & 5 & Construct & Timestamp, symbol, price, volume columns.\\
Eligible & 4 & Filter/update & Drop the zero price; calculate price times volume.\\
Enriched & 4 & Join & Add the matching venue; these keys do not multiply rows.\\
Totals & 2 & Aggregate & Accumulate volume and notional by symbol and venue.\\
Report & 2 & Update/order & Divide totals and explicitly sort by symbol.\\
\bottomrule
\end{tabularx}
\end{center}
The report contains AAPL with volume $30$, notional $3040$, and VWAP
$3040/30$; MSFT has volume $20$, notional $4030$, and VWAP $201.5$.
Keeping the division after aggregation expresses
$\sum_i p_i v_i / \sum_i v_i$, rather than an unweighted mean of prices.
The explicit final order makes this summary's presentation independent of
the first occurrence of each symbol in the ticks.

\paragraph{Observed execution boundary.}
Running the example with the existing release executable and
\code{--no-history} and \code{--report-planner} produced those two result rows and
reported statement execution because there were no lazy table sources to plan
against. Named intermediate results in this run are therefore separate
statement results. The sequence above describes data dependencies; it is
not a claim that the entire script became one fused physical pipeline.
The numeric output validates this small workload, not the parallel strategy
gates or compiled/interpreted parity. No engine rebuild or benchmark was
performed for this check.

Replacing the constructed input with a lazy file source would add demand
analysis and decoder behavior to the trace. Its exact physical plan should
be recorded from that run before making claims about pushdown or fusion.
Similarly, adding windows later will introduce persistent ordered state whose
boundaries need a separate treatment.
